
% 来源：EPUB part217.xhtml
\section{做自己喜欢做的事情}

\BookQuestion{网友：如何找到自己的激情所在？}

做自己喜欢做的事情，有点利润之上的追求。（2014-02-02）

% 来源：EPUB part218.xhtml
\subsection{得到你想要的，珍惜你得到的}

找工作最重要的是找自己喜欢做的工作。只有做自己喜欢做的事情，时间长了以后才最有机会做得很好。很多人找工作时都是首先衡量眼前的收入，换工作时也是这样，于是一辈子都可能在“循环以复”。（2014-06-06）

\BookQuestion{网友：就业选择是听长辈的还是听自己的？}

长辈的观点经常是对的，也经常不太对，我要总是按照长辈的观点去做的话，肯定不是今天这个样子。个人认为最重要的是你自己要有是非观，要有长远的看法。凡事如果你能想着10年20年，你大概就不会那么困惑了。（2019-04-01）

\BookQuestion{网友：诺贝尔物理学奖得主丁肇中说小孩子哪懂什么兴趣，只是历史比较好学，容易考得好，就觉得喜欢历史。我想想好像也有一定的道理，怎么样才能知道某个东西是兴趣而不仅仅是一个偷懒的理由呢？}

总有一天会知道的。家长可以帮忙，但最好不要强制。我猜大多数人可能一生都无法从事自己一开始就喜欢的行业，但有部分人可能后来会喜欢上自己一开始不喜欢的行业。（2010-07-05）

其实这也是我困惑的事情，为什么那么多人愿意长时间干着自己不喜欢的事情呢？我猜最重要的原因可能是这些人对自己没信心，我很难想到别的理由。

做对的事情，把事情做对——指的就是发现是错的事情就要改，不管多大的代价都会是最小的代价。

干着自己不喜欢的事情应该算是错的吧？我见过很多人一直就这么且待且抱怨着地过了很久很久。

我觉得人们选什么样的生活都有道理，但抱怨着却不打算改变就没啥道理了。（2014-09-21）

\BookQuestion{网友：好多事有兴趣但未必能做得成，没兴趣的事又已经能维持一个稳定的生活。}

这是个基本逻辑问题。没人说过干自己喜欢的事情就一定会成功的。做对的事情要成功也得有把事情做对的能力和资源。另外，做自己喜欢的事情至少在做的过程当中是快乐的，不然“喜欢”二字何来？（2014-09-24）

\BookQuestion{网友：大多数人的工作是为了生活，收入的高低是衡量工作满意度的重要标准之一，做自己喜欢的事情又有收入挺幸福的。}

区别一下：得到你想要的，珍惜你已经得到的（巴菲特语：Get what you want, want what you get）。琢磨一下很有意思。据说这是某种对成功和幸福的定义。（2014-05-26）

\BookQuestion{网友：什么是自己真正喜欢做的事情，这个问题需要经过仔细思考和反省才能找到答案么？还是应当很自然就知道的？目前个人在困惑的阶段，能找到的回答是先不做自己不喜欢的事情。}

是的，如果你正在做着自己非常不喜欢的事情，那就该想想了。其实很多人一辈子都找不到自己喜欢的事情做，其中相当多的人可能一直就没真的找过，因为一开始随便找了份工作，然后就年复一年地被这个工作拖进去了。不过，也有很多人喜欢自己的工作，这部分人往往就能做得更好些。（2020-10-15）

\BookQuestion{网友：对于没有天分轻松考取大学的孩子，寒窗苦读是做正确的事吗？还是对于穷人家的孩子这是正确的？}

找到自己喜欢做的事情应该是最重要的，和穷富无关。上好大学也许对找工作有帮助，但那个帮助区间其实是很短的，很快就没人在乎你是哪个大学毕业的了。

\BookQuestion{网友：很多年轻人也不容易找到自己喜欢的事，寒窗苦读考大学或许是发现并有效追求自己爱好的一个概率较大的途径吧。}

很容易就找到自己喜欢干的事情的人恐怕更少。多数人可能从来就没认真找过，但努力找赚钱多的工作单位的人还是不少的。一般来讲，做自己喜欢的工作的人收入可能会相对高些吧。（2018-10-09）

\BookQuestion{网友：看一本书谈到个观点，人们喜欢做的事情往往是自己做得比别人好的事情。觉得挺对的，就像读书时，往往学得好的课程越喜欢学一样。}

不知道作者有没有喜欢吃的东西，不知道怎么才能吃得比别人好。做得比别人好看起来是属于外在的东西。如果持这种观点的人，遇到比自己好的人的话就会马上放弃了。很多找不到自己喜欢的事情的人大多都没真认真找过什么是自己喜欢的事情，或者总是想着“钱多、活少、离家近”这类的东西。喜欢打高尔夫的人大多容易明白这个道理，因为绝大多数喜欢打高尔夫的人都不是因为自己比别人打得好。（2014-09-23）

\BookQuestion{网友：对拖延症有什么看法？}

我也有拖延症，好像没啥办法。据说和基因有关。

\BookQuestion{网友：所有拖延都源自不喜欢吗？}

这个说法可能有道理，喜欢的事情一般不会拖延。（2020-12-06）

\BookQuestion{网友：怎么看待马斯克说巴菲特的工作很无聊这件事？}

巴菲特很喜欢自己的工作，所以总是“跳着踢踏舞”去上班，他才不在乎别人说啥呢。（2020-10-15）

\BookQuestion{网友：这世界真有天才吗？还是天才就是学习方法和兴趣的加成？}

我不是那么相信天才的。我总觉得天才其实就是因为喜欢所以花的时间比别人多很多很多，于是就有了好很多很多的结果。你可以看看罗杰·费德勒在达特茅斯学院2024年毕业典礼上的致辞，那个讲话非常好。

\BookQuestion{网友：大道花在打高尔夫上的时间比花在投资上的时间多很多，然而在投资上获得的成就却又那么高，所以在我的心中大道就是天才啊。}

大道花在投资上的时间是很久以前花的，简单讲叫理解生意。就是说我花在努力理解生意的本质上的时间比99.9\%的人都要多很多很多。（2024-06-24）

\BookQuestion{网友：我从小就对企业家精神有浓厚的兴趣，并在初中开始接触投资市场。现在在美国读博。虽然在科研上过得去，但躁动的心却从未停止……您觉得年轻人应该去风口行业闯荡一下么？}

建议你看下巴菲特1998年在佛罗里达大学的那个演讲。可以多看几次，看到明白为止哈。这个演讲（交流）我听了不下10次，非常值得仔细听。（2020-12-03）

我其实想说的是：You are saving your sex for your old years（你在把性留到老年再享受）。（2020-12-04）

“我对大学和商学院的毕业生通常有两个建议。最重要的建议是到你欣赏/热爱的领域工作，听起来像你得喜欢你的经历，这很好。但是我们谈谈我的第二个建议，那就是坐上正确的火车，也就是说，朝着正确的方向前进。商学院没有所谓的‘搭上正确的火车’课程，但它确实重要。你可以是普通乘客，但如果你坐对了火车，它会载你走很长的路。管理你的职业生涯就像投资一样——困难的程度并不重要。因此，通过搭上正确的火车，你可以节省自己的金钱和痛苦。”（巴菲特）

很多人肯定很希望30年前能看懂这段话。当然现在看懂也不错，30年后就知道有多大差别了。

\BookQuestion{网友：怎样找到自己的“正确的火车”？}

如果你现在还不知道的话，30年后就一定会知道的，至少旁人会看出来。（2013-12-23）

\BookQuestion{网友：我有很多理想，但有一个共同的前提条件：等我有钱后。}

也许有点钱之上的追求会让你更有钱些。（2014-09-22）

\BookQuestion{网友：因为做得好所以喜欢什么事和喜欢做什么事然后把事情做好，虽然结果看起来差不多，因果关系却是反的。苹果不是因为有利润所以把产品做好，而是因为产品做得好所以才有利润。}

追着钱跑和让钱追着跑的区别。（2014-09-23）

“金钱是我把自己极其喜欢做的事情做到极致后所产生的一种附带结果。”（巴菲特）

前段时间在中欧我也说过钱是自己喜欢做的事情的副产品，不知道巴菲特有没有专利。（2012-11-16）

\BookQuestion{网友：找工作您有什么经验？如何分辨这家企业是否真的有利润之上的追求呢？}

我找工作也没啥经验。感觉是：只有当自己是个有利润之上的追求的人才会在乎公司是否有，在乎的人才会真的在意。（2013-03-08）

\BookQuestion{网友：怎样才能找到自己喜欢的事情？}

你必须去找才会知道怎么找。

如果没有喜欢的事情就确实很难找，如果有就不难。（2022-01-23）

\BookQuestion{网友：我老是对这个感兴趣对那个感兴趣，不能集中精力做一件事。请问有什么办法能抗干扰，聚焦做事吗？}

确定你真的对某件事情感兴趣吗？比较常见的事情是有些人只是对赚钱感兴趣，但不知道什么能赚钱，所以想去尝试各种听说赚钱更快的办法，不知道你是不是属于这种情况。（2019-08-08）

\BookQuestion{网友：价值投资的逻辑能否用在找工作上？比如你希望你服务的企业有好的商业模式和护城河，好的管理和企业文化。当然，如果你要投身到一家创业公司，盈利模式还没清楚的情况下，是不是比较难用以上逻辑判断呢？}

这个逻辑用在找工作上会非常好，很可能会改变一个人的一生。当然，开始找工作时往往人们没有那么多的时间和精力去了解得很透彻，但起码有这个逻辑的人会多少多花些时间在这方面考虑，而且一旦发现找错工作了就一定要尽快去找适合自己的工作。做对的事情往往体现在发现是错的事情就要尽快改。（2014-11-18）

无论如何，在找工作时一定不要仅仅是为了找一份有薪水的工作，一定要有些利润之上的追求。我当年去中山工作有一个很重要的原因是自己喜欢玩游戏，现在回想起来当时这个决定对我后来的影响是非常非常大的。（2014-06-15）

投身到创业公司的道理也一样，至少你想投身的这家创业公司一定要有什么东西特别触动你才行，要知道创业的过程是非常艰难的，没有利润外的追求很难走多远。我确实看到过一些人进入创业公司是因为创业公司会给期权，但那个东西对没价值的公司而言是完全没价值的。（2013-03-07）

% 来源：EPUB part219.xhtml
\subsection{越是迷惘的时候越是要往远处看}

越是迷惘的时候，越是要往远处看。习惯往远处看会相对容易做出正确的决定，长期而言大概率结果会比短视的人们好很多。（2022-07-02）

\BookQuestion{网友：最近一段时间常常会彷徨，怀疑自己是否有些傻。人究竟应该从哪个角度看待自己的人生？}

觉得会彷徨就说明你还不傻，呵呵，正常人都会偶尔有这种感觉的。至于为什么我就不知道了，不太清楚你说的角度指的是什么？我个人对人生的理解是不妨碍别人的前提下享受人生，要尽所能帮助别人但不要成为负担。不过这种问题大把人比我回答得好，你用谷歌或百度或搜狗查查说不定能找到你喜欢的答案。（2010-03-10）

\BookQuestion{网友：如何更好应对时间压力和生活压力？}

好像没什么办法。时间压力让你充实，生活压力让你有动机，都不是啥坏事啊。你要不太明白，反过来想想就明白了。（2010-03-28）

\BookQuestion{网友：大道说过大学的时候悟出快乐是过程中来的。有目标才有过程，想知道大道大学期间给自己定的目标是啥？工作之后呢？}

大学前阶段失去了目标，所以有段时间非常迷惘。后来发现不开心是因为失去了目标，所以开始认真考虑这个事情，但很长时期都没有答案。能从大学毕业已经是运气的事了。工作之后依然很长时间没有找到目标，直到开始经营小霸王后才开始觉得又找到自己喜欢做的事情了。

（2022-01-31）

\BookQuestion{网友：20岁的年纪应该给未来投资些什么？}

学会爱和健康比什么都重要。

\BookQuestion{网友：爱和健康，自认为那是有钱人的玩意呀！}

没有爱和健康的人其实也很难有钱，就算有钱了也不快乐。（2013-01-18）

% 来源：EPUB part220.xhtml
\subsection{为创业而创业的人多数很难成功}

\BookQuestion{网友：对于那些有创业理想的大学生，您有什么建议和忠告吗？第一，光有（创业）理想就去创业是不够的。}

第二，能否创业和是不是大学生好像并没有直接的联系。

一般而言（我个人观点，没验证过），学历越高，创业的难度越大，大概是机会成本的原因。

有很多专门的课讲如何创业，现在还有什么创业工厂也出来了。呵呵，讲这些课的老师大多是没有创过业的人，理论的东西讲讲可以，但到实际中还是要靠多磨练才能真正明白。如果创业工厂是由没创过业的人来搞，那就更有点高深了。

印象中有两种人创业成功的概率好像高一些：一种是无路可走的人，由于没有退路，他们会很拼命和专注，所以机会大些，比如像我。

另一种就是看到路的人，也就是创业开始时就有很强的理想和一定的想法（往往是得到了某种验证的想法），比如比尔·盖茨。

当然，也有为创业而创业的人成功的例子，比如惠普的两个创始人（《基业长青》里面有提到）。

总而言之，创业好像也没个公式，唯一可以肯定的就是创业失败的概率非常高。所以，除非你真是有可持续的高度热情（包括打不倒的精神）或者你真是无路可走了，否则还是多考虑考虑吧。（2010-03-20）

\BookQuestion{网友：可持续的高度热情一般人是做不到的。}

呵呵，所以创业成功的人很少！

\BookQuestion{网友：为啥说您没有退路呢？}

研究生毕业以后，找不到想走的路，就跑到广东去了。总的来讲，我真算是运气很好，最后找到了自己喜欢干的事情。（2010-03-21）

\BookQuestion{网友：那你是不是归为有能力而幸运的人？}

可以这么说。其实每个人都有自己的能力，对自己过于低估和高估都容易有问题。我比较客观看待自己。（2010-03-22）

\BookQuestion{网友：我在一家知名公司上班十多年了，这个行业我也挺喜欢，但是在公司上班却不能让我有成就感，感觉工作都是重复性的，我是否需要通过创业来实现人生的价值？}

创业的成功率很低，我没看过数据，但肯定比投资赚钱的比例还要低，应该至少在5\%以下（以活10年计算）。如果以成就感不够而创业的，估计成功率应该还要低很多。（2019-05-20）

如果你对创业有任何犹豫的话，可能就还没到时候。（2020-12-05）

创业的人是不会问别人是不是该去创业的。我认为能够创业成功的人都必须要有非常大的冲动（或者说自认为有很好的想法），到处问别人是不是该创业的人如果去创业多数是要失败的，或者说为创业而创业的人多数很难成功。（2010-09-06）

补充：我的意思其实不是反对创业，是说在创业前必须要有创业失败的精神准备。现在大环境下，只要真有好主意，创业的条件比我们那个时候好多了。我们那个时候都没听说过风险投资是啥。（2014-06-06）

到处问别人是不是该创业的人最好是别创业，失败的概率更高。（2025-01-09）

\BookQuestion{网友：关于创业时机，请谈谈体会。}

其实无论怎么准备，创业失败的概率都会远大于成功。看看那些成功的创业家都做过哪些准备吧，比如马云？如何创业可能是个巨大的话题，但个人体会，绝大多数创业成功者都是有追求的人，这种追求经常能克服创业过程中的许多巨大的困难。（2013-02-13）

每个创业成功的人可能经历都很不一样。我自己当年就是因为喜欢玩游戏，就一头扎进去了。（2020-11-30）

\BookQuestion{网友：对于初创公司，如果进入类似游戏这种充分竞争的战场，是不是兵家大忌？是不是也要像投资一样，要一直筛选，直到选到供需失衡（供小于需）的赛场？}

我觉得创业最重要的是创业者会觉得自己的产品或想法能够提供点别人提供不了的东西。如果只是想做生意，成功的概率会小很多。供需失衡的“赛场”大概没人可以在创业前就“选到”的，就算碰上了也很快会掉进红海。举个简单的例子：茅台的供不应求可不是一个创业者创业前计划好的。（2020-10-24）

\BookQuestion{网友：一直都想做点事情，但是一直都做不大，不知道自己能做些什么，觉得自己也没有什么真正擅长的东西。}

平平淡淡才是真，不是每个人都一定要当马云的。（2020-12-03）

\BookQuestion{网友：您认为一个初创企业最重要的是什么？}

我们当初想的都是活下来，理想都是后来才慢慢有的。现在情况不一样，信息很发达，人们可以看得远一些了。

\BookQuestion{网友：请问消费者导向是一开始就有还是后来有的？}

消费者导向和本分一样，必须是骨子里的东西，后天假装会是做不到的。（2024-10-09）

\BookQuestion{网友：我在实践中发现一个不可为的事项：那些一定要按甲方要求喝酒才能达成的业务，不可为。不管潜在业务量有多大，第一时间放弃，一定是最有利的选择。}

如果你确认对方是代表甲方逼你喝酒的话，我同意你的逻辑。早年我在小霸王时就碰到过这种情况。有个哈尔滨的客户在饭桌上对我说：干了这杯，不然这生意就不做了。当时我直接站起来客气滴说：那你们继续喝，我先走了。从此再无人逼我喝酒。甲乙双方是互惠关系，谁欺负谁都不对。

补充一下，那时候我确实不喝酒。现在喜欢喝点好红酒。（2021-09-25）
